
\section{\BookRecordIndex{T0003}$J$-全纯曲线的\BookRecordIndex{T0227}连通和}\label{sec:10-2}

本节将构造一族映射
\[
  \mathcal M(A^{0,\infty};J^{0,\infty})
  \longrightarrow C^\infty(S^2,M),
  \qquad
  (u^0,u^\infty)\longmapsto
  u^R:=u^0\mathbin{\#_R}u^\infty,
\]
它为每个 $J$-全纯球面的\BookRecordIndex{T0226}连通对 $(u^0,u^\infty)$ 指派一条近似 $J^R$-全纯球面 $u^R$。
更准确地说，该映射只定义在\BookRecordIndex{T0246}模空间 $\mathcal M(A^{0,\infty};J^{0,\infty})$ 的紧子集上，
而且对每个紧子集，只在 $R$ 充分大时定义。近似 $J^R$-全纯曲线
$u^R:S^2=\mathbb C\cup\{\infty\}\to M$ 在远离 $0$ 处与 $u^0$ 一致，在非常靠近
$0$ 处则与重标度曲线 $u^\infty(R^2z)$ 一致。条件
$u^0(0)=u^\infty(\infty)$ 保证这些映射在中心圆周 $|z|=1/R$ 上近似相等。我们将把这个
圆周加厚为宽度约为 $1/(\delta R)$ 的\BookRecordIndex{T0170}环域，并用\BookRecordIndex{T0188}截断函数构造\BookRecordIndex{T0254}粘合曲线 $u^R$，使它在该
\BookRecordIndex{T0170}环域上等于交点 $u^0(0)=u^\infty(\infty)$。在极限 $R\to\infty$ 时，我们得到\BookRecordIndex{T0139}格罗莫夫
紧性的逆过程：曲线 $u^R$ 将在定义~\ref{def:5-2-1} 的意义下收敛到对
$(u^0,u^\infty)$。

下面给出细节。固定常数 $c,\kappa>0$，并假设 $J^0,J^\infty$ 以这个常数 $\kappa$
满足~\eqref{eq:10-1-1}。取 $\varepsilon>0$ 小于 $M$ 关于由 $\omega$ 和
式~\eqref{eq:10-1-1} 中\BookRecordIndex{T0088}殆复结构 $J$ 经~\eqref{eq:2-1-1} 所确定度量的\BookRecordIndex{T0093}单射半径。固定光滑
\BookRecordIndex{T0188}截断函数 $\rho:\mathbb C\to[0,1]$，使得
\[
  \rho(z)=
  \begin{cases}
    0,& |z|\le1,\\
    1,& |z|\ge2.
  \end{cases}
\]

对每个 $(u^0,u^\infty)\in\mathcal M(c)$ 和任意满足 $0<\delta<1$ 且 $\delta R$
充分大的两个实数 $\delta,R$，构造近似 $J^R$-全纯曲线
\[
  u^{\delta,R}=:u^R:=u^0\mathbin{\#_R}u^\infty:S^2\longrightarrow M,
\]
它满足
\begin{equation}\label{eq:10-2-1}
  u^R(z)=
  \begin{cases}
    u^\infty(R^2z),& |z|\le\delta/(2R),\\
    u^0(0)=u^\infty(\infty),& \delta/R\le|z|\le1/(\delta R),\\
    u^0(z),& 2/(\delta R)\le|z|.
  \end{cases}
\end{equation}
因此，当 $z$ 非常接近 $0$ 时，$u^R$ 是 $u^\infty$ 的重标度；它在 $|z|=1/R$ 周围的
中心\BookRecordIndex{T0170}环域上为常值；而在该\BookRecordIndex{T0170}环域外足够远处等于 $u^0$。为在剩下的两个\BookRecordIndex{T0170}环域上定义 $u^R$，
把曲线 $u^0,u^\infty$ 的交点记作
\[
  x:=u^0(0)=u^\infty(\infty),
\]
并使用该点邻域中的指数映射。我们断言
\[
  |z|<\tan(\varepsilon/c),\quad |z|<\kappa
  \quad\Longrightarrow\quad
  d_J(u^0(z),x)<\varepsilon,
\]
以及
\[
  |z|>\cot(\varepsilon/c),\quad |z|>1/\kappa
  \quad\Longrightarrow\quad
  d_J(u^\infty(z),x)<\varepsilon.
\]
事实上，由~\eqref{eq:10-1-1}，当 $|z|<\kappa$ 时
$J_z^0=J_{1/\bar z}^\infty=J$；又有
$\|du^0\|_{L^\infty}\le c$ 和 $\|du^\infty\|_{L^\infty}\le c$（因而 $u^0,u^\infty$
都是 Lipschitz 常数为 $c$ 的连续映射）；而 $S^2$ 上的 Fubini--Study 度量满足
$d_{\mathrm{FS}}(z,0)=\tan^{-1}(|z|)$。

对 $|z|<\min\{\kappa,\tan(\varepsilon/c)\}$，令
$\zeta^0(z)\in T_xM$；对
$|z|>\max\{1/\kappa,\cot(\varepsilon/c)\}$，令
$\zeta^\infty(z)\in T_xM$；它们分别是唯一满足范数小于 $\varepsilon$ 且
\[
  u^0(z)=\exp_x(\zeta^0(z)),
  \qquad
  u^\infty(z)=\exp_x(\zeta^\infty(z))
\]
的切向量。现在对 $\delta/(2R)\le|z|\le2/(\delta R)$ 定义
\begin{equation}\label{eq:10-2-2}
  u^R(z):=u^{\delta,R}(z)
  :=\exp_x\left(
    \rho(\delta Rz)\zeta^0(z)
    +\rho\left(\frac{\delta}{Rz}\right)\zeta^\infty(R^2z)
  \right).
\end{equation}
只要
\begin{equation}\label{eq:10-2-3}
  0<\delta<1,
  \qquad
  \delta R>\frac{2}{\tan(\varepsilon/c)},
  \qquad
  \delta R>\frac{2}{\kappa},
\end{equation}
这一定义便是良好的。此外，式~\eqref{eq:10-2-2} 的第一项只在 $|z|\ge1/(\delta R)$
时非零，第二项只在 $|z|\le\delta/R$ 时非零。不过，映射 $u^R$ 并非
$J^R$-全纯曲线；我们将看到，其误差在 $R\to\infty$ 时趋于零。我们的目标是在 $u^R$
附近构造一条真正的 $J^R$-全纯曲线 $\widetilde u^R:=\widetilde u^{\delta,R}$。

\begin{remark}\label{rem:10-2-1}
固定常数 $0<\delta<1$。则曲线 $u^R:=u^{\delta,R}$ 在
定义~\ref{def:5-2-1} 的意义下收敛到对 $(u^0,u^\infty)$。更准确地说，$u^R(z)$ 在
$S^2\setminus\{0\}$ 的紧子集上一致连同各阶导数收敛到 $u^0(z)$，而
$u^R(z/R^2)$ 在 $S^2\setminus\{\infty\}$ 的紧子集上一致连同各阶导数收敛到
$u^\infty(z)$。下面构造的 $J^R$-全纯曲线 $\widetilde u^R$ 将以同样方式收敛到
$(u^0,u^\infty)$。
\end{remark}

\begin{remark}\label{rem:10-2-2}
考察一族在原\BookRecordIndex{T0226}连通对 $u^0,u^\infty$ 与其\BookRecordIndex{T0394}预粘合 $u^R$ 之间过渡的曲线对，会很有用。为此，
令 $r:=\delta R$，并考察由下式定义的曲线
$u^{0,r},u^{\infty,r}:S^2\to M$：
\[
  u^{0,r}(z):=
  \begin{cases}
    u^R(z),& |z|\ge1/r,\\
    u^0(0),& |z|\le1/r,
  \end{cases}
  \qquad
  u^{\infty,r}(z):=
  \begin{cases}
    u^R(z/R^2),& |z|\le r,\\
    u^\infty(\infty),& |z|\ge r.
  \end{cases}
\]
这样，$u^{0,r}$ 和 $u^{\infty,r}$ 分别是 $u^0$ 和 $u^\infty$ 的小扰动，并且分别在
$z=0$ 和 $z=\infty$ 附近被压平（即变成常值）。由下面的引理~\ref{lem:10-4-3}，当
$r\to\infty$ 时，$u^{0,r}$ 在 $W^{1,p}$ 范数下收敛到 $u^0$，而 $u^{\infty,r}$ 在
$W^{1,p}$ 范数下收敛到 $u^\infty$。因为 $du^0$ 和 $du^\infty$ 对所有
$(u^0,u^\infty)\in\mathcal M(c)$ 都一致地由 $c$ 控制，所以该收敛在所有这些对上一致。
虽然 $u^{0,r}$ 和 $u^{\infty,r}$ 不在 $C^1$ 范数下收敛，但这个 $W^{1,p}$ 收敛足以
说明：在用\BookRecordIndex{T0263}平行移动等同定义域空间与靶空间后，当 $r\to\infty$ 时，$D_{u^{0,r}}$ 和
$D_{u^{\infty,r}}$ 分别在算子范数下收敛到 $D_{u^0}$ 和 $D_{u^\infty}$。这由
第~\ref{sec:3-5} 节给出的 $D_u$ 公式立即得出。
\end{remark}

\begin{example}\label{ex:10-2-3}
考察 $M=S^2=\mathbb C\cup\{\infty\}$ 配备标准复结构的情形。全纯曲线
\[
  u^0(z)=1+z,
  \qquad
  u^\infty(z)=1+1/z
\]
满足所需的
$u^0(0)=u^\infty(\infty)=1$\sourcefootnote{原书此处误印为
$u(0)=v(\infty)=1$；本例只定义了 $u^0,u^\infty$，故按紧邻定义唯一补正其上标与字母。}。
上述映射 $u^R=u^{\delta,R}:S^2\to S^2$ 满足
\[
  u^R(z)=
  \begin{cases}
    1+1/(R^2z),& |z|\le\delta/(2R),\\
    1,& \delta/R\le|z|\le1/(\delta R),\\
    1+z,& 2/(\delta R)\le|z|.
  \end{cases}
\]
附近的 \BookRecordIndex{T0003}$J$-全纯曲线由
\[
  \widetilde u^R(z)
  =u^\infty(R^2z)+u^0(z)-1
  =\frac{z^2+z+1/R^2}{z}
\]
给出；它们如注~\ref{rem:10-2-1} 所述收敛到对 $(u^0,u^\infty)$。
\end{example}

\begin{remark}\label{rem:10-2-4}
应把\BookRecordIndex{T0254}粘合看作\BookRecordIndex{T0141}格罗莫夫收敛的逆过程。由于构造 \BookRecordIndex{T0003}$J$-全纯曲线很困难，一般过程颇为复杂。
不过，它仍适用于简单得多的\BookRecordIndex{T0339}稳定曲线情形，其中所有映射都是平凡的。例如，考察带四个
\BookRecordIndex{T0066}标记点的亏格零\BookRecordIndex{T0339}稳定曲线\BookRecordIndex{T0246}模空间 $\overline{\mathcal M}_{0,4}$。如第 D.7 节所述，通过\BookRecordIndex{T0185}交比
可将它与 $S^2$ 等同。开\BookRecordIndex{T0072}层 $S^2\setminus\{0,1,\infty\}$ 由元素
$\tau_z=[0,1,\infty,z]$ 组成，另有三个特殊元素 $\tau_0,\tau_1,\tau_\infty$，其定义域
各有两个分支。在这个背景下，\BookRecordIndex{T0141}格罗莫夫收敛描述当 $z\to a$（$a=0,1,\infty$）时元素
$\tau_z$ 如何收敛到 $\tau_a$；\BookRecordIndex{T0254}粘合则描述如何在
$\overline{\mathcal M}_{0,4}$ 中构造 $\tau_a$ 的一个邻域。

由于这个邻域有两个实维，其元素由一个复参数 $\epsilon$ 描述。采用当前记号并取 $a=0$，
可将 $\tau_0$ 等同于\BookRecordIndex{T0339}稳定曲线 $[z_1^0,z_2^0,z_3^0,z_4^0]$；它的定义域是两个
$S^2$ 副本 $S_0^2,S_\infty^2$ 的不交并，其中 $S_0^2$ 中的 $0$ 与
$S_\infty^2$ 中的 $\infty$ 等同，\BookRecordIndex{T0066}标记点为
\[
  z_2^0=1,
  \quad z_3^0=\infty\quad\text{位于 }S_0^2,
  \qquad
  z_1^0=0,
  \quad z_4^0=1\quad\text{位于 }S_\infty^2.
\]
于是，对给定参数 $|\epsilon|<1$，$\tau_\epsilon$ 是\BookRecordIndex{T0339}稳定曲线
$[z_1^\epsilon,z_2^\epsilon,z_3^\epsilon,z_4^\epsilon]$，其中
\[
  \begin{aligned}
  z_2^\epsilon=z_2^0=1,
  &\qquad z_3^\epsilon=z_3^0=\infty &&\text{于 }S_0^2,\\
  z_1^\epsilon=\epsilon z_1^0=0,
  &\qquad z_4^\epsilon=\epsilon z_4^0=\epsilon &&\text{于 }S_\infty^2.
  \end{aligned}
\]
因此 $\tau_\epsilon=[0,1,\infty,\epsilon]$。注意，当
$\epsilon=|\epsilon|$ 为实数时，在 $\epsilon=1/R^2$ 的对应下，这个稳定四元组可自然地
与\BookRecordIndex{T0394}预粘合 $u^R$ 的定义域等同。

我们将对实\BookRecordIndex{T0255}粘合参数完成主要\BookRecordIndex{T0254}粘合构造，但结果立即推广到复参数
$\epsilon=|\epsilon|e^{i\theta}$，因为乘以 $e^{i\theta}$ 是让两个球面彼此相对旋转的
等距映射。事实上，若族 $\{J_z^0,J_z^\infty\}$ 只依赖于 $|z|$，则整个构造都是
$S^1$-不变的。（见 McDuff \srcite{268} 以及第~\ref{sec:10-9} 节中定理~\ref{thm:10-9-1} 的证明。）
\end{remark}
